% Regression: the body restates one generated bond among many.  A lattice
% draws its bonds from the row and column policy, and before issue 5571 the
% body could not reach one of them: an authored wire over the same pair was
% recorded first and the generated bond inked on top of it, while the
% typed-port spelling of that pair raised TKZ-PORT-CONSUMED because the
% generated bond had already claimed both ports.  The frame populates and the
% body overrides -- the rule section 4 states for atoms -- so an authored
% index wire joining two adjacent cells IS that pair's bond, and the frame
% mints no second one under it.  The topology is stated once either way: the
% same contraction count, the same silhouettes, and the same exposed
% boundary, whether the restatement carries ink of its own or none.
\documentclass{standalone}
\usepackage{tenkz}
\begin{document}
\tenkzkernel
\tndeclare{species}{blocked}{hue=source:red}
% the control: every bond comes from the policy
\begin{tenkz}[rows={wire,wire,wire}, cols=3,
  west=open, east=open, north=open, south=open]
  \tn{} & \tn{} & \tn{}\\
  \tn{} & \tn{} & \tn{}\\
  \tn{} & \tn{} & \tn{}
\end{tenkz}
\qquad
% two of the twelve restated with nothing of their own: the same figure,
% one by its cells and one through its typed ports
\begin{tenkz}[rows={wire,wire,wire}, cols=3,
  west=open, east=open, north=open, south=open]
  \tn{} & \tn{} & \tn{}\\
  \tn{} & \tn{} & \tn[name=c]{}\\
  \tn{} & \tn{} & \tn[name=d]{}
  \tnwire{(1,1)}{(1,2)}
  \tnwire{c.270}{d.90}
\end{tenkz}
\qquad
% the same two carrying a species and a stroke: the edge under discussion
\begin{tenkz}[rows={wire,wire,wire}, cols=3,
  west=open, east=open, north=open, south=open]
  \tn{} & \tn{} & \tn{}\\
  \tn{} & \tn{} & \tn[name=e]{}\\
  \tn{} & \tn{} & \tn[name=f]{}
  \tnwire[species=blocked]{(1,1)}{(1,2)}
  \tnwire[species=blocked, stroke=dashed]{e.270}{f.90}
\end{tenkz}
\end{document}
